\section{已完成的主要研究工作}

\subsection{原理瓶颈与优化框架}

\begin{frame}{已完成的主要研究工作}{Two-tone编解码原理与性能瓶颈}

    \setbeamerfont{caption}{size=\scriptsize}
    \begin{columns}[T]
        \column{0.46\textwidth}
        \footnotesize
        \begin{block}{Two-tone 的基本特征}
            \footnotesize
            \begin{itemize}
                \setlength{\itemsep}{2pt}
                \item 编码仅由移位和异或构成，移位由地址偏移实现
                \item 解码可抽象为代入、消除和修复过程
                \item 核心瓶颈是\textbf{内存访问}，不是 XOR 算术
            \end{itemize}
        \end{block}

        \column{0.54\textwidth}
        \begin{figure}
            \centering
            \includegraphics[width=\textwidth,height=0.50\textheight,keepaspectratio]{Example of Two-tone Decoding.pdf}
            \vspace{-0.7em}
            \caption{Two-tone解码举例}
        \end{figure}
        \vspace{-0.4em}
        \footnotesize
        \begin{itemize}
            \setlength{\itemsep}{2pt}
            \item 冗余访存和时间局部性差
            \item 横向/纵向遍历难以同时适配所有场景
            \item 解码与修复分离导致额外遍历
        \end{itemize}
    \end{columns}

\end{frame}

\begin{frame}[shrink=10]{已完成的主要研究工作}{FastTT整体优化框架}

    \setbeamerfont{caption}{size=\scriptsize}
    \begin{columns}[T]
        \column{0.40\textwidth}
        \begin{block}{设计原则}
            \small 不改变 Two-tone 数学结构，围绕\textbf{内存访问路径}做系统级优化。
        \end{block}
        \vspace{0.2em}
        \footnotesize
        \begin{block}{四项协同优化}
            \footnotesize
            \begin{itemize}
                \setlength{\itemsep}{3pt}
                \item \textbf{EPM}：统一数据/校验恢复
                \item \textbf{CFWP}：窗口化提升缓存局部性
                \item \textbf{NMA}：相邻块合并降低写放大
                \item \textbf{TPS}：按丢失模式选择遍历路径
            \end{itemize}
        \end{block}
        \centering
        \footnotesize 编码侧：CFWP+NMA\quad 解码侧：四项协同

        \column{0.60\textwidth}
        \centering
        \includegraphics[width=0.88\textwidth,height=0.42\textheight,keepaspectratio]{Two-tone Optimization Overview.pdf}
        \vspace{0.2em}
        \scriptsize FastTT整体优化框架
    \end{columns}

\end{frame}

\subsection{四项优化}

\begin{frame}{已完成的主要研究工作}{EPM 与 CFWP：减少冗余遍历并提升局部性}

    \footnotesize
    \begin{columns}[T]
        \column{0.5\textwidth}
        \begin{block}{EPM：统一解码抽象}
            \footnotesize
            \begin{itemize}
                \setlength{\itemsep}{2pt}
                \item 将丢失数据块映射到存活校验块，丢失校验块映射到自身
                \item 把校验块修复融入解码流程，省去独立 repair 遍历
            \end{itemize}
        \end{block}
        \centering
        \includegraphics[width=0.84\textwidth,height=0.34\textheight,keepaspectratio]{Erasure-Parity Mapping Decoding Demonstration.pdf}
        \vspace{0.2em}
        \scriptsize 作用：减少冗余访存，统一处理所有丢失组合

        \column{0.5\textwidth}
        \begin{block}{CFWP：缓存友好窗口划分}
            \footnotesize
            \begin{itemize}
                \setlength{\itemsep}{2pt}
                \item 将块切分为窗口，窗口内完成相关移位异或
                \item 首次访问时初始化，避免全量预清零带来的额外开销
            \end{itemize}
        \end{block}
        \centering
        \includegraphics[width=0.92\textwidth,height=0.34\textheight,keepaspectratio]{Cache-Friendly Window Demonstration.pdf}
        \vspace{0.2em}
        \scriptsize 作用：在数据离开缓存前完成更多有效计算
    \end{columns}

\end{frame}

\begin{frame}{已完成的主要研究工作}{NMA 与 TPS：降低写放大并加速常见故障}

    \footnotesize
    \begin{columns}[T]
        \column{0.5\textwidth}
        \begin{block}{NMA：相邻合并访问}
            \footnotesize
            \begin{itemize}
                \setlength{\itemsep}{2pt}
                \item 相邻数据块先在 SIMD 寄存器内异或合并
                \item 合并后再写回校验块，减少校验块写入次数
            \end{itemize}
        \end{block}
        \centering
        \includegraphics[width=0.94\textwidth,height=0.36\textheight,keepaspectratio]{Neighbor-Merging Access Pattern.pdf}
        \vspace{0.2em}
        \scriptsize 作用：降低写放大，提升编码端吞吐

        \column{0.5\textwidth}
        \begin{block}{TPS：遍历模式选择}
            \footnotesize
            \begin{itemize}
                \setlength{\itemsep}{2pt}
                \item 单块丢失直接走纵向单遍快路径
                \item 避免完整代入/消除流程，每码字仅访问一次
            \end{itemize}
        \end{block}
        \centering
        \includegraphics[width=0.84\textwidth,height=0.36\textheight,keepaspectratio]{Vertical Traversal Pattern Demonstration.pdf}
        \vspace{0.2em}
        \scriptsize 作用：显著加速最常见的单块失效场景
    \end{columns}

\end{frame}
